\section{Related work}
\label{sec:relatedwork}

\subsection{Situational Impairments and Stress in Interactive Systems}
HCI research has traditionally evaluated interaction techniques under controlled laboratory conditions, often assuming stable cognitive and physical capabilities. Real-world interaction, however, frequently occurs in dynamic contexts that temporarily disrupt these capabilities. \citet{sears2003computers} introduced the concept of situationally induced impairments and disabilities (SIID) to describe such transient limitations in otherwise able-bodied users, arising from contextual factors such as environmental conditions, multitasking, and cognitive demands. Subsequent work has refined SIID as a broad class of context-dependent disruptions affecting interaction across diverse settings~\cite{li2026survey}.

This concept has recently been extended to immersive environments, where perception, action, and feedback are tightly coupled and interaction is therefore especially sensitive to contextual variation~\cite{li2026survey}. Walking and physical encumbrance reduce interaction accuracy and movement efficiency~\cite{li2025weight}, while ambient noise~\cite{li2025estimating} and lighting conditions~\cite{liu2025understanding} affect text entry, pointing performance, and task completion time. These findings indicate that behaviour observed under controlled conditions may not generalize to real-world immersive use.

Beyond environmental and physical factors, situational contexts can induce cognitive disruptions that substantially shape interaction. Acute stress, in particular, has been shown to alter attentional control~\cite{sarsenbayeva2019measuring}, cognitive load~\cite{hou2025cognitive}, and task strategies~\cite{mcguire2023stress}. Within HCI, \citet{sarsenbayeva2019measuring} demonstrated that acute stress affects mobile target acquisition, visual search, and text entry, reducing both accuracy and completion time. Complementary work in affective HCI and motor control suggests that emotional states can modulate low-level motor behaviour: affective priming influences pointing performance under Fitts' Law paradigms~\cite{joo2026quantifying}, and anxiety or cognitive stress can alter physiological tremor and fine motor variability~\cite{budini2022mental}. This line of work, however, focuses on affective priming or single-task motor performance rather than acute social stress in immersive multi-task settings. Taken together, prior work~\cite{sarsenbayeva2019measuring, joo2026quantifying, budini2022mental} establishes that situational and affective states can influence interaction, yet how acute stress manifests in immersive environments across tasks with different cognitive and perceptual demands remains unclear.

\subsection{Inducing and Measuring Stress in Immersive Environments}
A substantial body of research has examined the effects of acute stress on cognition, showing that it impairs working memory, reduces cognitive flexibility, and disrupts executive functioning~\cite{arnsten2009stress,starcke2016effects}, as well as higher-order processes such as decision making and attentional control~\cite{qin2009acute, starcke2012decision}. Thus, stress introduces systematic changes across cognitive functions that are essential for interactive behaviour.

The Trier Social Stress Test (TSST)~\cite{kirschbaum1993trier} is one of the most widely used paradigms for inducing acute stress in controlled settings. By combining public speaking and mental arithmetic calculations under social evaluation, it reliably elicits both physiological and psychological responses. In HCI, the TSST has been used to examine stress effects on mobile interaction, with measurable impacts on target acquisition, visual search, and text entry~\cite{sarsenbayeva2019measuring}, and on reliance decisions during interaction with automated systems~\cite{mcguire2023stress}.

Immersive environments have emerged as a complementary platform for studying stress and emotion. \citet{liszio2018relaxing, liszio2019interactive} showed that VR can both induce and regulate stress, with effects comparable to real-world TSST paradigms. Building on this, \citet{zhang2026seconds} designed a VR-based stress inoculation training system for novice physicians, identifying support strategies such as self-regulation aids, procedure guidance, and sensory support. While that work targets intervention design in a medical training scenario, our study examines how acute stress affects underlying VR interaction components, which can in turn inform when and what kinds of support are needed. Related VR affective computing research further shows that immersive scenes can elicit emotions~\cite{jiang2024immersive} and that interaction itself can modulate emotional experience~\cite{kuang2026understanding}, supporting VR as a controlled yet ecologically relevant platform for studying how stress- and emotion-related state changes unfold during interaction.

Stress is typically assessed through a combination of subjective, physiological, and behavioural measures, including self-report scales such as the State-Trait Anxiety Inventory (STAI)~\cite{spielberger1971state}, heart rate and heart-rate variability~\cite{kim2018stress}, and behavioural indicators such as completion time and error rate~\cite{li2026survey,sarsenbayeva2019measuring}. Together, these measures capture both internal state changes and their observable interaction consequences. Most existing studies, however, focus on a single task context or non-immersive settings, leaving open how stress-induced state changes translate across different interaction modalities within the same immersive environment. 

\subsection{Task-Based Probes and Multiple Resource Theory}
Understanding how stress affects interaction requires tasks that engage different combinations of cognitive, perceptual, and response resources. Multiple Resource Theory (MRT)~\cite{wickens2008multiple} provides a widely used framework for modelling human performance, proposing that interference and workload depend on the degree to which tasks draw on overlapping resource dimensions. MRT organises resources along three primary dimensions: processing stages (perception, cognition, and responding), perceptual modalities (primarily visual and auditory), and processing codes (spatial and verbal); Wickens' 2008 formulation~\cite{wickens2008multiple} further nests focal and ambient visual processing within the visual modality. MRT has therefore been used extensively to reason about why tasks with different resource-demand profiles exhibit different levels of interference, workload, and performance vulnerability. This perspective is well suited to our study because target acquisition task, reading aloud task, and Tower of London task all occur within the same embodied VR setting, yet differ in their dominant resource demands: visuomotor control, verbal delivery, and executive spatial planning, respectively.

In HCI, task-based evaluation is commonly used to probe interaction performance across levels of complexity, including low-level perceptual--motor tasks such as target acquisition~\cite{fitts1964information,li2025estimating,li2025weight}, language-related tasks such as text entry~\cite{mackenzie2003phrase,sarsenbayeva2019measuring} and reading aloud~\cite{dillon2021real}, and higher-level cognitive tasks involving planning and problem solving~\cite{shallice1982specific,banakou2018virtually}. These tasks serve as controlled probes for isolating specific aspects of user performance. Recent work has further shown that different interaction tasks impose distinct cognitive demands that can be characterized using multimodal signals in immersive settings~\cite{hou2025cognitive}. However, such approaches primarily aim to identify task differences or estimate cognitive load, rather than to examine how a situational factor such as stress differentially affects performance across tasks. Prior work on stress and interaction, in turn, has typically treated performance as a single outcome, implicitly assuming uniform degradation under stress.

In this work, we adopt MRT as a task-selection framework to compare interaction performance across tasks with distinct resource profiles. Rather than directly measuring resource competition, we use task-based probes to examine how acute stress affects performance across perceptual -- motor, verbal, and higher-level cognitive domains. This multi-task perspective allows us to move beyond uniform-degradation assumptions and systematically investigate whether stress effects are task-dependent in immersive interaction contexts.

